\section{犯罪构成 - 定罪方法：三段论推理}
\begin{enumerate}
  \item 大前提是法律规定的一个罪的\zd{构成要件}（唯一标准），小前提是案件事实 \yjt 逻辑推导 \yjt 结论。
  \item \bs{构成要件要素}。构成要件的最小单位。其实就是一个个词语。分类。
        \begin{enumerate}
          \item 记述、\zx{规范}（不是指法律规范）。事实判断/\zx{价值评价}。\al{毒品的认定、淫秽物品的认定}。
          \item 客观、主观。行为外在的 \yjt 人、事、物、时、地/\zx{行为人内心}的 \yjt 故意、过失、目的、动机。
          \item 积极、消极。正面证明/\zx{反面否定犯罪成立}。\al{没有获得不正当利益的，不是行贿 \yjt 唯一个消极要素}。
          \item 成文、\zx{不成文}。明文规定/\zx{没有规定，但实质上必须具备}的要素。
        \end{enumerate}
  \item 案件事实的认定由刑事诉讼法负责，无法认定，作\zx{存疑时有利于被告}原则。
  \item 正确推导注意事项。
        \begin{enumerate}
          \item 避免颠倒大小前提。
          \item \zx{循环往复推导}，一行为触犯多个罪。
          \item 道德不能成为定罪的标准。\al{大义灭亲}。
        \end{enumerate}
\end{enumerate}
